% ICLR methods · 第二轮 Opus 5（飞书完整现在稿 + 构建手册 + tool 设计口径 + 作者本轮意见）
% 尚未替换 chapters/03_MobileGym++_20260915.tex
% 机械校正：节标题、We built、七领域、原图表；删掉模型捏造的对比表和 3.1 里的 52 tools。
% 3.5 只换开篇，Task design / Scale / Protocol 仍用飞书现稿。

\section{MobileGym++}
\label{sec:mobilegym++}
\label{sec:method}

\subsection{A gap in controllable commercial applications}
\label{sec:mgpp-env}

Mobile agents are increasingly moving toward real-world workflows across commercial applications, such as travel booking on Trip and online shopping on Taobao, where tasks are often long-horizon and grounded in real user needs. However, training and evaluation environments for such commercial applications remain largely missing, leaving agents with limited opportunities to practice and be evaluated in realistic mobile-use scenarios.

Real-device or Android-emulator setups face login, reset, and real-world side effects that make large-scale training impractical. Prior simulated-app work typically covers only a small slice of commercial apps and rarely provides resettable business state~\citep{rawles2025androidworld,kong2025mobileworld}.

We built MobileGym++, a browser-hosted mobile environment that carries commercial workflows in controllable simulated applications. The runtime extends MobileGym~\citep{wu2026mobilegym}, preserving its advantages in reset and programmatic evaluation while substantially expanding domain coverage, the number and richness of applications, and support for realistic cross-app workflows. Observations remain screen pixels, there is no accessibility tree, and the agent cannot read business state; device state can be exported and restored, and many instances can run in parallel. In this round we newly built 35 commercial apps, bringing the total with the predecessor to 64 ready-to-use applications spanning travel, shopping, finance, content, social, office, and local services. Each application maintains persistent business states, such as orders, budgets, and itineraries, which evolve through interaction and enable task verification based on the resulting state.

Figure~\ref{fig:mgpp-overview} summarizes the environment at a glance: commercial-app coverage, executable single- and cross-app tasks, the two interaction channels, and their shared terminal check.

\begin{figure}[t]
\centering
\includegraphics[width=\linewidth]{figures/mobilegym_fig1_20260913.pdf}
\caption{MobileGym++ combines broad commercial-app coverage with executable single- and cross-app tasks and two interchangeable interaction channels. GUI-only and Hybrid share the same initial state and terminal check.}
\label{fig:mgpp-overview}
\end{figure}

\subsection{35 high-fidelity commercial applications}
\label{sec:mgpp-apps}

We implemented 35 commercial clients as replicas of the corresponding real apps, each built by a human annotator paired with a coding agent, reproducing a business loop rather than a visual skin. The annotator records screen recordings and long screenshots of the core flows and key pages, then freezes the object model and its state transitions before any page is written: what an order, a ledger entry, or an itinerary is, and how each is created, modified, and closed. The coding agent implements pages and business functions around those objects; humans compare the replica against the real device, and mismatches go back into the annotator--agent loop. The procedure is specified in Appendix~\ref{app:mgpp-build}. Goods and orders, expenses and budgets, itineraries and changes are maintained as application state, rendered by the UI, and read at task termination.

Application count alone is not informative: a client carries a multi-step commercial task---locate, fill, confirm---only if it is structurally thick. On the 35 new clients the median navigation declaration is 37 routes, 175 interactive actions, and 88 transitions, and the median interaction depth from the root is 3, so a typical shortest GUI path already crosses several layers and the step lower bound comes from page structure rather than from the length of the instruction. As a depth reference, the thickest application in MobileGym has 100 routes and 265 actions and its deepest terminal takes 7 steps to reach; the new commercial clients sit inside that structural range rather than below it.

Table~\ref{tab:mgpp-compare} compares coverage (left) with interface and checking (right). It is not a score table. Per-app statistics are in Appendix~\ref{app:mgpp-apps}.

\begin{table}[t]
\caption{MobileGym++ beside existing controllable environments, on coverage and interface/checking rather than scores (task definitions and platforms preclude a shared protocol). Platforms: Emu.\ = emulator with real apps; Sim.\ = browser-hosted simulation; API\ = no phone GUI. Interaction: G\ = GUI; A\ = API; H\ = in-app tools. $\checkmark$ yes, $\sim$ partial, --- no commensurate statistic. Task counts keep each paper's unit and are not a difficulty ranking.}
\label{tab:mgpp-compare}
\begin{center}
\scriptsize
\setlength{\tabcolsep}{3.5pt}
\begin{tabular}{llrclcc}
\hline
 & \multicolumn{3}{c}{Coverage} & \multicolumn{3}{c}{Interface / check} \\
\cline{2-4} \cline{5-7}
Env. & Plat. & Apps & Tasks & Act. & X-app & Check reads \\
\hline
\textbf{MobileGym++ (ours)} & Sim. & 64 (+35) & 215 & G+H & $\checkmark$ & app business state \\
MobileGym & Sim. & 28 & 416 tmpl. & G & $\sim$ & device-state diff \\
AndroidWorld & Emu. & $\sim$20 & 116 tmpl. & G & $\sim$ & system state \\
AppWorld & API & 9 & 750 / 250 sc. & A & $\checkmark$ & app DBs \\
MobileWorld & Emu. & 20 & 201 & G+MCP & $\checkmark$ & mixed \\
\hline
\end{tabular}
\end{center}
\end{table}

\subsection{Realistic cross-application workflows}
\label{sec:mgpp-crossapp}

Several commercial applications sit on one resettable device, each keeping its own live business state: orders, expenses, itineraries, contacts. These objects change with on-screen actions and leave a trace in device state. Clipboard, system broadcasts, and inter-app navigation connect them, so a workflow can leave the current application, enter a second one, and rewrite that application's own business objects: a payment in shopping produces a matching expense in bookkeeping; a train number selected in the calendar enters an itinerary app as a change; an address from contacts enters a checkout flow. The second application receives a real business write, not a forwarded message.

Real commercial applications do behave this way on a phone, but not controllably: login, anti-bot defenses, and irreversible side effects make the same multi-application starting point impossible to restore. Here the whole device is one snapshot, so restoring it lets the same cross-application write be run again.

\subsection{Hybrid GUI and app-native tool use}
\label{sec:mgpp-tools}

GUI-only benchmarks assume an application never exposes a programmatic interface, and programmatic-only environments with no phone GUI assume the opposite. Real applications sit in between: they expose lookup, filtering, and drafting over their own records. MobileGym++ puts both paths in one world---the same task, the same initial state, and the same application-owned store.

Approximately 52 applications expose app-native tools. Tools are designed per application from that application's business objects---searching orders, filtering ledger entries, retrieving itinerary fields---rather than exported from the widget tree. A tool call occupies the same action slot as a tap, swipe, or text input while that application is in the foreground. Both channels write the same store, and every official task can be completed with GUI actions alone: Hybrid is a second path through the same world, not a requirement.

\begin{figure}[t]
\centering
\includegraphics[width=\linewidth]{figures/mobilegym_fig3_20260913.pdf}
\caption{Hybrid execution in MobileGym++. GUI actions and app-native tools operate on the same business state; a trajectory may hand off between channels before the checker reads the resulting state diff.}
\label{fig:mgpp-hybrid}
\end{figure}

\subsection{MobileGym++ Bench}
\label{sec:mgpp-bench}

Existing mobile-agent benchmarks mainly score single-app GUI sequences whose terminal is system state or a UI-level field~\citep{rawles2025androidworld,kong2025mobileworld}. Two properties are not given jointly: the terminal is a commercial write in an application's own store---an order, an expense, an itinerary---and the value of that write lives in another application's live state, so the second application is actually rewritten rather than notified. MobileGym++ Bench asks whether an agent can carry those writes to completion inside resettable simulated commercial applications. Task length comes from locating, filtering, confirming, paying, and writing back across applications.

\paragraph{Task design.}
Tasks are organized around executable user goals and checkable terminals. Core tasks close with a business operation such as placing an order, posting an expense, creating an itinerary, sending a message, or modifying an app object. The result of that operation lands in the application's own business state. A single-app task reaches that write only after locating and filling fields across deeper pages. Instructions use user-goal language and do not name widget labels, tool namespaces, or checker fields.

In a cross-application task, the value needed for a write is present in another app's live state rather than written as a fixed constant in the instruction. The agent must identify the source, carry the value across the app boundary, and write it into the correct field; the checker verifies the relation between the resulting business fields.

\paragraph{Scale and composition.}
The current set contains 215 tasks across seven business domains: 109 mention one app, 74 two, 27 three, and 5 four. Tasks are distributed over multiple applications and business objects, so evaluation is not dominated by a handful of apps.

\paragraph{Evaluation protocol.}
Each task has three parts: a natural-language user goal, a bound initial state, and a programmatic terminal check. Success is decided by a programmatic checker: it reads the before/after diff of application business state, does not read the trace or screenshots, and does not use a model judge; different routes that accomplish the same write are equivalent. Allowed state paths are declared, and other mutations are recorded separately from success. The primary metric is task success rate. GUI-only and Hybrid are evaluated on the same tasks, initial states, and terminal checks. Mean steps on successful trajectories measure interaction efficiency conditional on completion; because failed trajectories often reach the step cap, we also report the success rate and the fraction reaching the cap to separate efficiency from reliability.
